% \newpage
% \begin{center}
% \normalfont\Large\bfseries MPC exercise: week two
% \end{center}
\setcounter{section}{0}
% \section{Two player LQR}
% Recall the one-player LQR formulation from the lecture 
% \begin{align*}
%         \min_u \quad &\sum_{k = 0}^{N-1} x_k^{\top}Qx_k + u_k^{\top}Ru_k + V_N(x_N),\\
%     \text{s.t.}\quad &x_{k+1} =Ax_k + Bu_k,
% \end{align*}
% where the agent has control over $u_k$. Assuming the value function is $V_k(x) = x^{\top}P_kx$, we obtain the LQR controller
% \begin{align*}
%     u_k &= -(R+B^{\top}P_{k+1}B)^{-1}B^{T}P_{k+1}Ax_k.
% \end{align*}
% Now consider a two-player scenario where the system has control inputs $u$ and $v$ with a shared state $x$, $x_{k+1} =Ax_k + B_1u_k + B_2v_k$. Agent one has access to control input $u$ while agent two has access to control input $v$. For agent one, the cost function for the controller is 
% \begin{align*}
%     \min_u \quad &\sum_{k = 0}^N x_k^{\top}Q_1x_k + u_k^{\top}R_1u_k + V_{N,1}(x_N),\\
%     \text{s.t.}\quad &x_{k+1} =Ax_k + B_1u_k + B_2v_k.
% \end{align*}
% For agent two, the cost function is 
% \begin{align*}
%     \min_v \quad &\sum_{k = 0}^N x_k^{\top}Q_2x_k + v_k^{\top}R_2v_k + V_{N,2}(x_N),\\
%     \text{s.t.}\quad &x_{k+1} =Ax_k + B_1u_k + B_2v_k.
% \end{align*}
% Notice that two cost functions are coupled via the shared states $x$. In this exercise, we will develop the best response strategy controller for each player (also called Nash Equilibrium in game theory, not required by this course or the final exam).
% \begin{enumerate}
%     \item \label{ex2:1} Assume the value function at stage $k+1$ is $V_{k+1,1}(x) = x_{k+1}^{\top}P_{k+1,1}x_{k+1}$ for agent one, and $V_{k+1,2}(x) = x_{k+1}^{\top}P_{k+1,2}x_{k+1}$ for agent two. Write down the cost function $f_1$ and $f_2$ for dynamic programming as a function of $x_k,u_k,v_k$ at stage $k$ for each agent.\vspace{5\baselineskip}
%     \item From \ref{ex2:1}, the cost function at stage $k$ is $f_1$ for agent one and $f_2$ for agent two. The best strategy is represented by the point where $\frac{\partial f_1(x_k,u_k,v_k)}{\partial u_k} = 0,\frac{\partial f_2(x_k,u_k,v_k)}{\partial v_k} = 0$. This will give us the best response strategy controller, $\begin{bmatrix}
%         u_k\\
%         v_k
%     \end{bmatrix} = \begin{bmatrix}
%         *&*\\
%         *&*
%     \end{bmatrix}^{-1}\begin{bmatrix}
%         *\\
%         *
%     \end{bmatrix}x_k$. Complete those two matrices.\vspace{5\baselineskip}
% \end{enumerate}
\section{Understanding LQR}
Consider the following LQR formulation
\begin{align*}
    \min_u\quad &\sum_{k=0}^{N-1}x_{k}^{\top}Qx_{k} + u_{k}^{\top}Ru_k + x_{N}^{\top}Px_N,\\
    \text{s.t.}\quad &x_{k+1} = Ax_k + Bu_k, \forall k \in [N-1].
\end{align*}
\begin{enumerate}
    \item Explain how the controller will behave when\\
    1) $Q = 0, P = 0$;\\
    2) $R = 0$.
    
    \solution{If $Q=0$, $P=0$, the cost function becomes $\min_u\quad \sum_{k=0}^{N-1} u_{k}^{\top}Ru_k$, which is only related to the inputs. The optimal solution to this problem is $u_k = 0, \forall k \in [N-1]$.}
    
    \solution{If $R=0$, the objective becomes $\min_u\quad \sum_{k=0}^{N-1}x_{k}^{\top}Qx_{k} + x_{N}^{\top}Px_N$, which is not related to the control inputs $u$. The system will settle to $x = 0$ as fast as possible.}
    \item Explain the difference between finite horizon LQR and infinite horizon LQR ($N = \infty$). Why is the infinite horizon controller time invariant?
    
    \solution{For finite horizon LQR, the controller is $u_k = -K_k x_k$, where $K_k$ is obtained by solving the discrete Riccati equation. For infinite horizon LQR, $\lim_{k \rightarrow \infty}K_k = K$. The controller becomes $u_k = -Kx_k$ with $K$ being a constant matrix. Hence, the infinite horizon controller is time-invariant.}
\end{enumerate}
\section{Dynamic programming}
Consider the following optimal control problem
\begin{align*}
    V^*(x_0) = \min_u \quad &\gamma x_N^2 + \sum_{k = 0}^{N-1}u_k^2,\\
    \text{s.t.}\quad&x_{k+1} = x_k+ u_k,
\end{align*}
where only the cost of the final state is penalized with $\gamma \geq 0$. Your goal is to solve this problem with dynamic programming.
\begin{enumerate}
    \item Assume that $V_{k+1}(x_{k+1}) = hx_{k+1}^2$ and derive $V_k(x_k)$.
    
    \solution{The Bellman recursion is 
    \begin{equation*}
        V_k(x_k) = \min_{u_k} u_k^2 + V_{k+1}(x_{k+1}) = \min_{u_k} u_k^2 + h(x_k+ u_k)^2.
    \end{equation*}
    Taking the derivative of $V_k(x_k)$ with respect to $u_k$, we have 
    \begin{equation*}
        u_k + h(x_k + u_k) = 0,
    \end{equation*}
    and 
    \begin{equation*}
        u_k = -\frac{h}{1+h}x_k.
    \end{equation*}
    Plugging $u_k$ back into $V_{k}(x_k)$, we obtain
    \begin{align*}
        V_k(x_k) &= \frac{h^2}{(1+h)^2}x_k^2 + \frac{h}{(1+h)^2}x_k^2\\
         &=\frac{h}{1+h}x_k^2.
    \end{align*}
    }
    \item What is $V_N(x_N)$?
    
    \solution{$V_N(x_N) = \gamma x_N^2$}
    
    \item Compute the optimal value function $V^*(x_0)$ and the control law $\kappa^*(x_0)$ for a finite horizon $N$.

    \solution{Starting from $V_N(x_N)$, we know 
    \begin{align*}
        V_{N-1}(x_{N-1}) &= \frac{\gamma}{\gamma + 1}x_{N-1}^2,\\
        V_{N-2}(x_{N-2}) &= \frac{\gamma}{2\gamma + 1}x_{N-2}^2,\\
        V_{N-3}(x_{N-3}) &= \frac{\gamma}{3\gamma + 1}x_{N-3}^2.
    \end{align*}
    Hence, $V^*(x_0) = \frac{\gamma}{N\gamma + 1}x_{0}^2$. The optimal control law $\kappa^*(x_0) = -\frac{h_1}{h_1 + 1}x_0 = -\frac{\gamma/((N-1)\gamma +1)}{\gamma/((N-1)\gamma +1) + 1}x_0 = -\frac{\gamma}{N\gamma + 1}x_0$.}
    \item Compute the closed-loop pole of your system. Will your system converge to zero for all finite $N$ and $\gamma$?

    \solution{The closed-loop system is
    \begin{equation*}
        x_{k+1} = x_k + \kappa(x_k) = \frac{(N-1)\gamma + 1}{N\gamma + 1}x_k.
    \end{equation*}
    Hence, the closed-loop pole is $\frac{(N-1)\gamma + 1}{N\gamma + 1}$. The system will converge to zero for all finite $N$ and $\gamma$ since the pole is smaller than one.
    }
    \item What is your control law when $N\rightarrow \infty$? Will your system converge to zero for infinite $N$ and finite $\gamma$?

    \solution{$\lim_{N\rightarrow\infty}\kappa(x_k) = 0$. That means the control inputs will always be zero. The pole of the closed-loop system converges to one. Hence, the system will not move.}
\end{enumerate}